% chktex-file 0
\documentclass{ctexart}
\usepackage{graphicx} % Required for inserting images
\usepackage[margin=2.5cm]{geometry}
\usepackage{booktabs}
\usepackage{float}
\usepackage{soulutf8}
\usepackage{enumitem}
\usepackage{pifont}
\usepackage{pgfplots}
\usepackage{longtable}
\pgfplotsset{compat=newest}
\usepackage{tikz}
\usepackage{xcolor} % 用于颜色设置
\usepackage{listings} % 用于代码块
\usepackage{fontspec} % 关键：用于使用系统字体 (XeLaTeX/LuaLaTeX)
% Linux 环境中没有 Windows 宋体时，使用可用的中文衬线字体作为宋体替代
\setCJKmainfont{Noto Serif CJK SC}
\usepackage{authblk}
\usepackage[
    colorlinks=true,      % 使用颜色而非方框
    linkcolor=blue,       % 内部链接颜色
    citecolor=green,      % 引用链接颜色
    urlcolor=red,         % URL 链接颜色
    bookmarks=true,       % 生成 PDF 书签
    bookmarksopen=true,   % 默认展开书签
    pdfstartview=FitH     % 页面宽度适应窗口
]{hyperref}

\sethlcolor{yellow}%设置所需高亮颜色
\usepackage[most]{tcolorbox}
\tcbuselibrary{listings,skins,breakable}
% 页眉页脚
\usepackage{fancyhdr}
\pagestyle{fancy}

% 定义草绿色
\definecolor{grassgreen}{RGB}{124,179,66}
% 补全缺失的浅蓝色定义（防止编译报错）
\definecolor{lightblue}{RGB}{30,144,255} 

% 清空默认
\fancyhf{}

% 页眉
\fancyhead[L]{\textcolor{lightblue}{\kaishu ICS}}
\fancyhead[R]{\textcolor{lightblue}{\textit{2026.9-2027.1}}}

% 页脚页码居中
\fancyfoot[C]{\textcolor{lightblue}{\thepage}}

% 页眉、页脚线颜色
\renewcommand{\headrulewidth}{1pt}
\renewcommand{\footrulewidth}{1pt}
\renewcommand{\headrule}{\color{blue!30!white}\hrule height 1pt width\headwidth}
\renewcommand{\footrule}{\color{blue!30!white}\hrule height 1pt width\headwidth}
\newtcolorbox{fancybox}[2][blue]{ % 可选颜色参数
  enhanced,
  breakable,
    colback=#1!3!white,
    colframe=#1!18!white,
    boxrule=0.8pt,
    arc=6pt,
    outer arc=6pt,
  left=10pt,right=10pt,top=12pt,bottom=10pt,
  title={\textbf{#2}},
    fonttitle=\bfseries\color{#1!65!black},
  attach boxed title to top left={yshift=-2mm,xshift=5mm},
  boxed title style={
        colback=#1!8!white,
        colframe=#1!18!white,
        boxrule=0.6pt,
        arc=4pt,
  },
    shadow={1.5pt}{-1.5pt}{0pt}{black!10},
}

% 定义一套适合C++的配色方案
\usepackage{tabularx}
\newfontfamily{\consolasfont}{Ubuntu}

\newcommand{\bluecode}[1]{{\color{blue}\consolasfont #1}}
\lstdefinestyle{MyCppStyle}{
    language = C++,
    % 确保这里使用 \consolasfont
    basicstyle = \consolasfont, 
    keywordstyle = \color{blue!70!black}, 
    commentstyle = \color{green!50!black}, 
    stringstyle = \color{red!60!black}, 
    numberstyle = \tiny\color{gray}, 
    numbers = left, 
    frame = single, 
    rulecolor = \color{lightgray}, 
    breaklines = true, 
    backgroundcolor = \color{white}, 
    tabsize = 4, 
    showstringspaces = false, 
    xleftmargin = 2em, 
    escapeinside=``,
}
% --- lstdefinestyle HTML 风格定义 ---
\lstdefinestyle{mhtml}{
    language = HTML,
    basicstyle = \consolasfont\small, 
    tagstyle = \color{blue!70!black}\bfseries,
    keywordstyle = \color{purple!80!black},
    commentstyle = \color{green!50!black}\itshape, 
    stringstyle = \color{red!60!black}, 
    numberstyle = \tiny\color{gray}, 
    numbers = left, 
    frame = single, 
    rulecolor = \color{lightgray}, 
    breaklines = true, 
    backgroundcolor = \color{white}, 
    tabsize = 4, 
    showstringspaces = false, 
    xleftmargin = 2em, 
    escapeinside=``,
    morekeywords={als, doctype, html, head, title, meta, link, body, div, span, h1, h2, h3, h4, h5, h6, p, a, img, ul, ol, li, table, tr, th, td, pre, code, script, style},
    morekeywords=[2]{class, id, href, src, rel, type, charset, lang, name, content, style},
    keywordstyle=[2]\color{purple!80!black},
}

% ==================== 汇编代码环境 ====================
% 与 MyCppStyle 共用同一套字体、配色、行号与边框，只换语言和背景色
\definecolor{asmbg}{RGB}{237,246,255}% 很浅的蓝，等于 lightblue!8!white

\lstdefinelanguage{x86asm}{
    sensitive    = false,       % 大小写不敏感，%RAX 与 %rax 都能识别
    alsoletter   = {.},         % 让 .globl / .quad 这类伪指令整体参与匹配
    morecomment  = [l]{\#},     % AT&T 语法的注释符
    morecomment  = [l]{;},      % Intel 语法的注释符
    morecomment  = [l]{//},
    morecomment  = [s]{/*}{*/},
    morestring   = [b]{"},
    morestring   = [b]{'},
    morekeywords = {%            指令与伪指令，沿用 C 的关键字颜色
        mov, movb, movw, movl, movq, movabs,
        movzbl, movzwl, movslq, movsbl, movswl,
        lea, leaq, push, pushq, pop, popq, xchg,
        add, addq, sub, subq, imul, imulq, mul, idiv, idivq, div,
        inc, dec, neg, not, and, andq, or, orq, xor, xorq,
        sal, sar, shl, shr, cmp, cmpq, test, testq,
        jmp, je, jz, jne, jnz, jg, jge, jl, jle, ja, jae, jb, jbe, js, jns,
        sete, setne, setg, setge, setl, setle, seta, setb,
        cmove, cmovne, cmovg, cmovge, cmovl, cmovle,
        call, callq, ret, retq, leave, nop, nopw, syscall, endbr64,
        .globl, .global, .text, .data, .bss, .rodata, .section,
        .align, .p2align, .byte, .word, .long, .quad, .string, .asciz,
        .zero, .size, .type, .file, .ident
    },
    morekeywords = [2]{%         寄存器
        rax, rbx, rcx, rdx, rsi, rdi, rbp, rsp, rip,
        r8, r9, r10, r11, r12, r13, r14, r15,
        eax, ebx, ecx, edx, esi, edi, ebp, esp,
        ax, bx, cx, dx, si, di, bp, sp,
        al, bl, cl, dl, sil, dil, bpl, spl, ah, bh, ch, dh,
        r8d, r9d, r10d, r11d, r12d, r13d, r14d, r15d,
        xmm0, xmm1, xmm2, xmm3, ymm0, ymm1
    }
}

\lstdefinestyle{MyAsmStyle}{
    language         = x86asm,
    basicstyle       = \consolasfont,
    keywordstyle     = \color{blue!70!black},
    keywordstyle     = [2]\color{purple!80!black},
    commentstyle     = \color{green!50!black},
    stringstyle      = \color{red!60!black},
    numberstyle      = \tiny\color{gray},
    numbers          = left,
    frame            = single,
    rulecolor        = \color{lightgray},
    breaklines       = true,
    backgroundcolor  = \color{asmbg},
    tabsize          = 4,
    showstringspaces = false,
    xleftmargin      = 2em,
    escapeinside     = ``,
}

% 直接可用的环境：\begin{asmcode} ... \end{asmcode}
\lstnewenvironment{asmcode}[1][]{\lstset{style=MyAsmStyle, #1}}{}

\usepackage{hyperref}
\hypersetup{
    colorlinks=true,      
    linkcolor=blue,       
    filecolor=magenta,    
    urlcolor=cyan,        
    pdftitle={我的文档标题}, 
}

% ==================== 正文部分 ====================
\begin{document}

% 开启无页眉页脚模式（针对封面和目录）
\pagestyle{empty}

% ------------------ 美化封面页 ------------------
\begin{titlepage}
    \centering
    % 顶部 TikZ 几何几何装饰条（全宽 bleed 效果）
    \begin{tikzpicture}[remember picture, overlay]
        \fill[blue!50!white] (current page.north west) rectangle ([yshift=-2cm]current page.north east);
        \fill[grassgreen] (current page.north west) ++(0,-2cm) rectangle ([yshift=-2.2cm]current page.north east);
    \end{tikzpicture}
    
    \vspace*{3.5cm}
    
    % 大标题
    {\fontsize{30pt}{36pt}\selectfont\bfseries\color{blue!80!black} 计算机系统导论（ICS）}\\[0.6em]
    % 副标题或英文标识
    {\fontsize{14pt}{18pt}\selectfont\color{gray}\textsf{Introduction to Computer System}}\\[1cm]
    
    % 居中装饰线
    \centering\makebox[4cm]{\color{grassgreen}\rule{5cm}{1.5pt}}
    
    \vfill
    
    % 使用 tcolorbox 包装的精致作者信息框
    \begin{tcolorbox}[
        enhanced,
        width=0.55\textwidth,
        colback=blue!5!white,
        colframe=blue!5!white,
        boxrule=1.2pt,
        arc=6pt,
        sharp corners=downhill, % 别致的对角切角效果
        shadow={2pt}{-2pt}{0pt}{black!20},
        halign=center,
        valign=center,
        top=15pt, bottom=15pt
    ]
        {\large\bfseries\color{black!80} Author：} {\large\kaishu Simson} \\[0.6em]
        {\large\bfseries\color{black!80} Date：} {\large\kaishu September 2026}
    \end{tcolorbox}
    
    \vfill
    \vspace*{2cm}
    
    % 底部标识
    {\color{gray!80}\small\the\year\ \copyright\ All Rights Reserved.}
\end{titlepage}

\pagestyle{fancy}

\newpage
\tableofcontents
\newpage

% 设置从当前页（正文第一页）开始采用阿拉伯数字编码，并自动重置为第 1 页
\pagenumbering{arabic}

\section{基本介绍}

本课程非常著名，使用的教材是CSAPP《深入理解计算机系统》

\subsection{ICS的主要内容}
\begin{itemize}
    \item 机器语言和编译（编译优化生成机器语言）
    \item 程序性能评估和优化
    \item 存储结构组织和管理
    \item 网络技术和协议
    \item 并行计算相关知识
\end{itemize}

\subsection{CSAPP的主要章节}
\begin{enumerate}
    \item 数据类型及其表示：了解计算规律和数学算术规律的差别
    \item 汇编程序和程序的机器级表示
    \begin{itemize}
        \item 了解程序执行模型，查找相关的程序错误（底层汇编代码工具）
        \item 程序性能调优
        \item 软件开发和系统软件
        \item 安全性分析（汇编代码静态分析）
    \end{itemize}
    \item 处理器体系结构
    \item 存储器层次结构：程序和数据在栈上的分布
    \item 异常控制流（定制shell程序实验）
    \begin{itemize}
        \item 进程控制
        \item 信号和信号处理
    \end{itemize}
    \item 虚拟内存：数据的布局和数据组织
    \item 算法性能分析
    \begin{itemize}
        \item 内存性能影响程序性能
        \item 程序编写符合缓存大小和性能
    \end{itemize}
    \item 计算机网络和并发编程（代理服务器）
\end{enumerate}

\section{信息的表示和处理}

\subsection{信息的二进制表示}

Reason : Binary presentation is the most practical system to use.

\subsubsection{信息量的单位}

bits是最基本的单位。

Byte = 8 bits

\subsubsection{Boolean Algebra}
逻辑运算的代数表示，与，或，非，异或。

用boolean代数的表示方法，可以将与或非异或运算对应到交并补差等集合运算

\subsubsection{数据的表示}

\paragraph{Unsigned 数据}

直接转换二进制存储

\[
B2U(X) = \sum_{i = 0}^{w - 1} 2 ^ i 
\]

范围：

\[
Umin = 0 ,\qquad
Umax = 2^{w} - 1
\]
\paragraph{Signed 数据}

负数原来使用源码表示，正负0不好区分等问题最终弃用，改为补码（Two's Complement）。
最高位为符号位，0为正数，1为负数

\[
B2T(X) = -2 ^ {w - 1} + \sum_{i = 0}^{w - 2} 2 ^ i 
\]

负数的补码等于模减去该负数的绝对值；对于某一确定的模，某数减去小于模的另一个数，总可以用另一个数字加上该数字的补码来确定

范围：

\[
Tmin = -2^{w-1} , \qquad
Tmax = 2 ^{w-1} - 1
\]

\begin{fancybox}{Example}
\[10 \to 01010\]
\[-10 \to 10110 \]
\[-1 \to 11111111...\](全是1)
\end{fancybox}

在编程中，有不同计算机的long类型有关键字可以跨平台进行。

\begin{lstlisting}[style=MyCppStyle]
#include<limits.h>

LONG_MAX
LONG_MIN

\end{lstlisting}

在C语言中，常量如果不加U,则默认为有符号数。加U则认为无符号数。

类型的转换相当于先转换为二进制表示，再对于同一个二进制串使用不同的方式进行解读。
在进行数值比较的时候，如果有无符号数同时出现，自动转换为无符号数字。

存储大小如果变大，在高位复制填入符号位直到填满为止，操作称为扩充。
存储大小如果缩小，直接按照所需的位数进行溢出截断，相当于对$2^n$进行取模运算

\subsection{浮点数}

浮点数是大模型数值表示的物理基础，大模型的很多工程技巧本质上都是与浮点数的局限性做折中权衡。训练和推理大量使用低精度浮点数，大模型训练几乎都使用混合参数。大模型部署的时候经常做量化，比如转换FP16为INT8或者INT4，虽然目标格式是整数，但是整个过程依赖于对浮点数值分布的理解。

\subsubsection{二进制分数Fractional Binary Numbers}

\paragraph{二进制分数数学表示}

二进制分数的基本形式：$1011.11_2$，二进制分数和二进制整数的表示遵从同样的规律：

\[
x = \sum_{n = -j}^{i} 2^n 
\]

{\sloppy 分数的乘除法和整数一样，都是靠小数点移位来实现的。需要注意一个数值的特殊表示，即 $0.1111111111\ldots_2$，它总是小于 $1$，记为 $1 - \epsilon$。\par}

二进制分数具有一定的范围局限性：只能精确表示小数部分可以被拆分为0.5的幂次之和的。其他的有理数会变成二进制循环小数。同时由于数据类型的大小限制，只能表示一定范围和一定精度的浮点数。

\paragraph{浮点数表示的IEEE标准}

浮点数标准来源于IEEE Standard754,在此之前有很多个人的独立表示。IEEE754标准适合于所有主流CPU，不同数据类型的转换过程中会存在很多错误（譬如Ariana的坠机来源于将一个64位浮点数转换为16位带符号整数时，产生了溢出异常）。下面对浮点数的表示技巧进行讲解。

浮点数表示方法来源于二进制科学记数法，以一个数字为例：
\[
1.10110011001012_2 \times 2^13
\]
根据二进制科学计数法表示，可以将二进制小数拆分为三个部分表示。
\[
(-1)^s M 2^E
\]

s为符号表示，M为有效位，E为指数位置。编码的时候按照数据类型的具体大小分配到三部分进行编码。
\begin{figure}[H]
\centering
\begin{tikzpicture}[font=\footnotesize]

  % ==================== 单精度 float：1 + 8 + 23 = 32 位 ====================
  \def\yS{2.2}   % 单精度条的下边界
  \def\bh{0.8}   % 条高

  \filldraw[fill=red!35,   draw=red!75!black,   line width=0.6pt] (0,\yS)    rectangle (0.19,\yS+\bh);
  \filldraw[fill=blue!25,  draw=blue!65!black,  line width=0.6pt] (0.19,\yS) rectangle (1.71,\yS+\bh);
  \filldraw[fill=green!30, draw=green!55!black, line width=0.6pt] (1.71,\yS) rectangle (6.08,\yS+\bh);

  \node at (0.95,\yS+0.4)  {阶码 8 位};
  \node at (3.895,\yS+0.4) {尾数 23 位};

  \node[anchor=south west] at (0,\yS+\bh+0.1) {\textbf{单精度浮点数 float}（共 32 位）};
  \node[anchor=east] (sgnS) at (-0.12,\yS+0.4) {符号位 1 位};
  \draw[->,red!75!black,line width=0.6pt] (sgnS.east) -- (0,\yS+0.4);

  % ==================== 双精度 double：1 + 11 + 52 = 64 位 ====================
  \def\yD{0}     % 双精度条的下边界

  \filldraw[fill=red!35,   draw=red!75!black,   line width=0.6pt] (0,\yD)    rectangle (0.19,\yD+\bh);
  \filldraw[fill=blue!25,  draw=blue!65!black,  line width=0.6pt] (0.19,\yD) rectangle (2.28,\yD+\bh);
  \filldraw[fill=green!30, draw=green!55!black, line width=0.6pt] (2.28,\yD) rectangle (12.16,\yD+\bh);

  \node at (1.235,\yD+0.4) {阶码 11 位};
  \node at (7.22,\yD+0.4)  {尾数 52 位};

  \node[anchor=south west] at (0,\yD+\bh+0.1) {\textbf{双精度浮点数 double}（共 64 位）};
  \node[anchor=east] (sgnD) at (-0.12,\yD+0.4) {符号位 1 位};
  \draw[->,red!75!black,line width=0.6pt] (sgnD.east) -- (0,\yD+0.4);

\end{tikzpicture}
\caption{单精度与双精度浮点数的存储结构（两行按同一比例绘制，条形长度正比于位数，故双精度条长为单精度的两倍）}
\end{figure}

\paragraph{三类位模式}

位域划分定下来之后，三段具体取什么值，由阶码 $exp$ 决定。所有位模式分为三类：

\begin{center}
\small
\begin{tabular}{llll}
\toprule
\textbf{阶码 exp} & \textbf{尾数 frac} & \textbf{类别} & \textbf{数值} \\
\midrule
全 0               & 全 0   & $\pm 0$                & $\pm 0$ \\
全 0               & 非全 0 & 非规格化数（Denormalized） & $(-1)^s \times 0.f \times 2^{1-bias}$ \\
非全 0 且非全 1    & 任意   & 规格化数（Normalized）     & $(-1)^s \times 1.f \times 2^{\,exp-bias}$ \\
全 1               & 全 0   & $\pm\infty$            & $\pm\infty$ \\
全 1               & 非全 0 & NaN                    & 非数（Not a Number） \\
\bottomrule
\end{tabular}
\end{center}

\textbf{规格化（Normalized）数}是最常见的一类，三个字段这样解释：

\begin{itemize}
    \item $E = exp - bias$，其中偏置值 $bias = 2^{k-1} - 1$，$k$ 为 $exp$ 字段的位数。单精度 $k = 8$，$bias = 127$；双精度 $k = 11$，$bias = 1023$。
    \item $M = 1 + f$，$f$ 是 $frac$ 字段的位模式所表示的小数部分。也就是说尾数最高位那个 1 是\textbf{隐含}的，不占存储空间。
\end{itemize}

偏置值取 $2^{k-1} - 1$，是为了让阶码的取值范围大致对称地跨在 0 两侧；同时这使得按无符号整数去比较 $exp$ 字段，就等价于比较两个正浮点数的大小，硬件比较电路因此可以简化。

\textbf{非规格化（Denormalized）数}：$exp$ 全 0 时，规定 $E = 1 - bias$（而不是 $0 - bias$，这样才能与规格化数的边界平滑衔接），同时 $M = f$，\textbf{不再有隐含的最高位 1}。
\begin{itemize}
    \item $frac$ 全 0：表示数值 0，但符号位仍然有效，因此有 $+0$ 与 $-0$ 之分（见下）。
    \item $frac$ 非全 0：表示最接近 0 的那些数。单精度最小的非规格化数是 $2^{-149}$。
\end{itemize}

\textbf{特殊值}：$exp$ 全 1 时，若 $frac$ 全 0 则表示无穷 $(\pm\infty)$，若 $frac$ 非全 0 则表示 NaN（Not a Number，非数）。

\paragraph{正零与负零}

$exp$ 全 0 且 $frac$ 全 0 时表示数值 0，但符号位并不因此失效，于是同一个数值 0 对应两个不同的位模式：IEEE 754 规定 $+0 = -0$，C 语言中 \texttt{0.0 == -0.0} 求值为真。所以\textbf{不能用 \texttt{==} 来区分正负零}，要判断符号得用 \texttt{signbit()}，或者去看 $1/x$ 的符号。\textbf{参与运算时二者不等价。}符号位在算术中照常起作用：

\begin{itemize}
    \item 除法：$1/(+0) = +\infty$，而 $1/(-0) = -\infty$。
    \item 乘法：结果的符号由两个操作数的符号\textbf{异或}决定，所以 $(-3.0) \times (+0) = -0$，而 $(-0) \times (-0) = +0$。
    \item 加法：在就近舍入模式下 $(-0) + (+0) = +0$。
    \item 取平方根：$\sqrt{-0} = -0$。
\end{itemize}

保留-0目的是让符号信息在运算中连续传播，不至于中途丢失。在数值分析和复变函数里，$-0$ 表示"从负方向趋近于 0"这一极限方向：$x \to 0^-$ 时 $1/x \to -\infty$，$\log(-0) = -\infty$。如果只保留 $+0$，左右极限就无法区分。
整数采用补码，不存在负零。整数因此弃用了符号-幅值表示。

C 语言中注意区分两个写法：\texttt{-0.0} 是合法的浮点字面量，得到负零；而 \texttt{-0} 是整数常量，它的值就是 0，不存在负零。此外 \texttt{printf("\%f", -0.0)} 会输出 \texttt{-0.000000}，负零的符号是可以被打印出来的。
\subsection{C语言的运算}

C语言中，将数据看作二进制向量进行处理。

\begin{itemize}
    \item 算术运算
    \item 逻辑运算
    \item 移位运算
\end{itemize}

\subsubsection{算术运算}

和位相关，直接按位进行Boolean代数运算。

\subsubsection{逻辑运算}

主要是逻辑与，逻辑或

\subsubsection{移位运算}

左移：右侧补0
右移：分算术右移和逻辑右移，逻辑右移直接向右移动高位补0,算术右移右移补符号位

\subsubsection{无符号数运算}

\paragraph{无符号加法}
首先是无符号数字的加法，首先是二进制加法，有可能发生溢出，相当于取模运算，最多取模一次

\[
Uadd(u, v) = u + v mod 2^w
\]

有符号数字的加法在bit位上有与无符号数完全相同的行为。如果计算结果超出了范围，可能发生正溢出，也可能发生负溢出。

\[
\boxed{
\begin{aligned}
\operatorname{TAdd}_w(u,v)
&=
\begin{cases}
u+v+2^w, & u+v<TMin_w \quad \text{\textbf{(NegOver)}} \\[4pt]
u+v, & TMin_w\le u+v\le TMax_w \\[4pt]
u+v-2^w, & TMax_w<u+v \quad \text{\textbf{(PosOver)}}
\end{cases}
\end{aligned}
}
\]

\paragraph{无符号数乘法}

乘法发生溢出的情况可能更多，可能需要算法包进行扩展，增大存储范围。

\[
Umult(u, v) = u \cdot v \bmod 2^w
\]

由于可能发生正负溢出，所以TMult有时候和UMult相比会出现不同的正负性。

乘法的具体操作是使用左右移动的方法，将乘数分解为二进制。左移k位，相当乘2的k次幂。

\paragraph{整数除法}

除法相对于乘法，采用右移的方式来实现。除法涉及到向上取整的问题。

\subsubsection{浮点数运算}
基本的思路是，首先精确计算结果，然后根据格式做变换和解释。可能会根据计算情况进行溢出截断和舍入操作。

\[
x +_f y = Round(x + y)
\]

\[
x \times_f y = Round(x + y)
\]

\paragraph{舍入规则}浮点数舍入规则为向偶数舍入。第一规则是向最近的数字靠拢，如果是中间数值距离是一样的，要向偶数舍入。其他的方法尤其是在求和的情况下有舍入的偏差，应当确认向最低的偶数位进行舍入，保证大规模运算的统计均值与真实数值一致。

计算机在执行舍入的时候，会设置在最低位进行按规则舍入，对于形如BBGRXXX，G为哨兵位即最低有效位；R为Round位具体执行舍入。Sticky bit在后面是要被扔掉的数字，用Sticky bits判断应该朝哪边进位。这一过程成为Postnormalize，在浮点数解释完成之后处理。

\paragraph{乘法规则}

浮点数乘法规则非常简单，符号位异或，指数位相加，小数frac位相乘。

如果乘法结果大于2,则指数数位+1，如果E超出了范围，报溢出错误。M应当舍入到合适的精度。

\paragraph{加法规则}

浮点数加法规则难于乘法，由于加法应当满足阶数对齐规则。不妨假设参与运算的浮点数满足$E_1 > E_2$。
通过移位运算（阶数小的数字向右移），将两个数的阶对齐。对齐之后应该移位就移位，该舍入就舍入，该报错就报错。

\begin{figure}[H]
\centering
\begin{tikzpicture}[font=\footnotesize, x=0.62cm]
  % ===== 公共位权轴 =====
  \foreach \i/\wt in {0/{2^{3}},1/{2^{2}},2/{2^{1}},3/{2^{0}},4/{2^{-1}},5/{2^{-2}},6/{2^{-3}},7/{2^{-4}}}
    \node[gray] at (\i+0.5,3.02) {$\wt$};
  \draw[gray,->] (-0.05,2.80) -- (8.05,2.80);
  \node[gray,anchor=west] at (8.15,2.80) {位权};

  % ===== 条形图一：x（阶 E_1 = 3），不移位 =====
  \foreach \i/\dig in {0/1,1/1,2/0,3/0,4/1,5/0,6/0,7/0}{%
    \filldraw[fill=blue!22,draw=blue!65!black,line width=0.5pt] (\i,1.90) rectangle (\i+1,2.60);
    \node at (\i+0.5,2.25) {\dig};}
  \node[anchor=east,blue!65!black] at (-0.15,2.25) {$x = 1.1001_2 \times 2^{3}$};

  % ===== 条形图二：y（阶 E_2 = -1），右移 4 位后 =====
  \foreach \i in {0,1,2,3}{%
    \draw[fill=white,draw=gray,dashed,line width=0.5pt] (\i,0.50) rectangle (\i+1,1.20);
    \node[gray] at (\i+0.5,0.85) {0};}
  \foreach \i/\dig in {4/1,5/0,6/0,7/0}{%
    \filldraw[fill=green!28,draw=green!55!black,line width=0.5pt] (\i,0.50) rectangle (\i+1,1.20);
    \node at (\i+0.5,0.85) {\dig};}
  \node[anchor=east,green!45!black] at (-0.15,0.85) {$y = 1.0_2 \times 2^{-1}$};

  % ===== 移位标注 =====
  \draw[<->,gray,line width=0.5pt] (0,0.28) -- (4,0.28);
  \node[gray,anchor=north] at (2,0.20) {右移 $\Delta E = E_1 - E_2 = 4$ 位，空位补 $0$};
\end{tikzpicture}
\caption{浮点数加法的对阶过程（以 $12.5 + 0.5$ 为例）。两行条形图共用同一条位权轴：$y$ 的阶小，尾数右移 $\Delta E = 4$ 位（虚线格为补入的 $0$），使两个数落在同一组位权上，再按列相加得 $1.1010_2 \times 2^{3} = 13$}
\end{figure}

\paragraph{运算性质}

浮点数加法和Abel群对比，对加法封闭，但是不支持结合律。除了Nan值和无穷值，所有元素都有加法逆元，除了特殊值满足不等式的加法性质。

浮点数和交换环相比，乘法满足交换性，由于溢出不满足结合率，同样不满足分配率，具有幺元1,不满足单调性，除了特殊值满足不等式的乘法性质。

\paragraph{类型转换}

将浮点数转换为整数：向0舍入，如果是未定义怎么转换的数值默认转换为Tmin。float有自身的转换规则，所以转换过程中会有舍入损失。
\subsection{信息的表示和存储}

内存就是一个大型的数组，以word为单位存储。程序通过地址访问内存中的程序。
内存中最基本的单位是字，一个地址占据的空间大小称为字长；这个参数决定了地址在内存中的分布。
系统会按照字节进行编址和寻址。常见的系统是32-bit或者64-bit，指的是一个地址空间对应的比特位。

一个多字节的数据，例如一个整数，占据4个字节。根据多字节数据在内存中排布的字节顺序，决定了两种存储模式。
大端法指的是：最小的有效位置放在最大的地址，反之小端法最小的有效数字放在最小地址（地址的标示是整数）， 
需要注意的是Internet的存储方式是使用大端法。

每个字符只占据一个字节，所以字符串的存储不受到大小端法的影响，可以用字符串起到标示和检查的效果。

\section{机器级编程:Machine Programming}

\subsection{处理器系统}

\subsubsection{Intel处理器}
Intelx86处理器系统是当前主流的处理器系统。Intel处理器的设计是不断进化的，Intel指令集性能较好，功耗较高。

关于Intel的指令集非常有名，简称CISC，叫作复杂指令集。指令集的内容很多且格式复杂，但是只有一小部分是在linux系统会见到的。复杂指令集的功耗让很难match到精简指令集系统的表现。

\subsubsection{x86 CLones： Advanced Micro Devices（AMD）}

AMD从前一直追随Intel架构，用微弱的性能换极低的能耗。在进化到64位系统之后首先主导开发Intel X86-64处理器。

\subsubsection{相关定义}

\paragraph{Architecture}

也称ISA，指指令集架构。指令集架构指的是在编写机器语言和汇编语言的时候需要遵循的指令集规范。（架构决定了如何解释机器代码和汇编代码）

\paragraph{MicroArcitecture}

cache sizes ，core frequency一类计算机内部的具体执行方式的总称，决定了程序具体被如何执行。

\subsubsection{机器语言和汇编代码视角下的计算机系统}

CPU处理器模型：
\begin{itemize}
    \item PC：Programm counter
    \item Register file
    \item Condition codes(条件码)
\end{itemize}

内存Memory模型：

\begin{itemize}
    \item Byte addressable array
    \item Code and user data
    \item Stack to support procedures
\end{itemize}

\begin{figure}[H]
\centering
\begin{tikzpicture}[font=\footnotesize, x=1cm, y=1cm,
    cpu/.style={draw=blue!65!black, fill=blue!8, line width=0.8pt, rounded corners=3pt},
    mem/.style={draw=green!55!black, fill=green!8, line width=0.8pt, rounded corners=3pt}]

  % ===== CPU：PC / Register file / Condition codes =====
  \draw[cpu] (0,0) rectangle (4.2,2.8);
  \node[blue!65!black] at (2.1,2.45) {\textbf{CPU}};
  \node[align=center] at (2.1,1.28) {PC (Program counter)\\[4pt] Register file\\[4pt] Condition codes};

  % ===== Memory：按字节编址的数组 =====
  \draw[mem] (9.2,0) rectangle (13.6,2.8);
  \node[green!45!black] at (11.4,2.45) {\textbf{Memory 内存}};
  \node[align=center] at (11.4,1.20) {Byte addressable array\\[4pt] Code and user data\\[4pt] Stack to support procedures};

  % ===== CPU $\to$ Memory：地址 =====
  \draw[->,blue!65!black,line width=1pt] (4.2,2.05) -- (9.2,2.05);
  \node[fill=white,inner sep=2pt] at (6.7,2.05) {地址 Address};

  % ===== Memory $\to$ CPU：指令 =====
  \draw[->,green!45!black,line width=1pt] (9.2,1.35) -- (4.2,1.35);
  \node[fill=white,inner sep=2pt] at (6.7,1.35) {指令 Instruction};

  % ===== CPU $\leftrightarrow$ Memory：数据 =====
  \draw[<->,red!70!black,line width=1pt] (4.2,0.65) -- (9.2,0.65);
  \node[fill=white,inner sep=2pt] at (6.7,0.65) {数据 Data};
\end{tikzpicture}
\caption{CPU 与内存之间的数据流动。PC 给出地址，内存按字节编址取出该地址上的指令或数据送回 CPU；load/store 经 Register file 读写数据，Condition codes 记录运算结果的状态。}
\end{figure}

\subsection{C程序的编译过程和反汇编}

\subsection{汇编语言}

汇编代码不是凭空写出来的，它由 C 代码经一条固定的流水线生成，也可以反过来从二进制里还原出来。这两个方向各有一套工具，是读懂机器级程序的基本功。

\begin{figure}[H]
\centering
\begin{tikzpicture}[font=\footnotesize, x=1cm, y=1cm,
    txtf/.style={draw=blue!65!black, fill=blue!8, line width=0.8pt, rounded corners=3pt,
                 align=center, minimum width=2.3cm, minimum height=1.5cm},
    binf/.style={draw=green!55!black, fill=green!8, line width=0.8pt, rounded corners=3pt,
                 align=center, minimum width=2.3cm, minimum height=1.5cm}]

  % ===== 五个中间产物：前三个是 ASCII 文本，后两个是二进制 =====
  \node[txtf] (f1) at ( 1.15,0.75) {sum.c\\[2pt] C 源文件};
  \node[txtf] (f2) at ( 4.40,0.75) {sum.i\\[2pt] 预处理后};
  \node[txtf] (f3) at ( 7.65,0.75) {sum.s\\[2pt] 汇编代码};
  \node[binf] (f4) at (10.90,0.75) {sum.o\\[2pt] 可重定位目标};
  \node[binf] (f5) at (14.15,0.75) {a.out\\[2pt] 可执行文件};

  % ===== 四步工具 =====
  \draw[->,line width=0.8pt] (f1) -- node[above,inner sep=1pt] {cpp} (f2);
  \draw[->,line width=0.8pt] (f2) -- node[above,inner sep=1pt] {cc1} (f3);
  \draw[->,line width=0.8pt] (f3) -- node[above,inner sep=1pt] {as}  (f4);
  \draw[->,line width=0.8pt] (f4) -- node[above,inner sep=1pt] {ld}  (f5);
\end{tikzpicture}
\caption{从 C 源文件到可执行文件的生成流程。cpp、cc1、as、ld 依次是预处理器、编译器、汇编器和链接器。前三步的产物（.c、.i、.s）都是 ASCII 文本，人可以直接读；后两步的产物（.o 和可执行文件）是二进制，只能用工具解析。gcc 的 \texttt{-E}、\texttt{-S}、\texttt{-c} 三个选项分别让它停在对应的阶段。}
\end{figure}

\subsubsection{从 C 到汇编：编译流水线}

以一段求数组和的 C 程序为例，它也是后面反复要用的例子：

\begin{lstlisting}[style=MyCppStyle]
#include <stdio.h>

long sum(long *a, long n) {
    long s = 0;
    for (long i = 0; i < n; i++)
        s += a[i];
    return s;
}

int main(void) {
    long a[] = {1, 2, 3, 4, 5};
    printf("%ld\n", sum(a, 5));
    return 0;
}
\end{lstlisting}

gcc 并不是一步把 .c 变成可执行文件的，它按上图把四个阶段依次串起来。用 \texttt{-E}、\texttt{-S}、\texttt{-c} 可以让它在任意一步停下，看清每个中间产物：

\begin{lstlisting}[style=MyCppStyle, language=sh]
$ gcc -E sum.c -o sum.i     # 预处理器 cpp：展开 #include 和宏
$ gcc -S sum.c              # 编译器 cc1：把 C 翻译成 sum.s
$ gcc -c sum.c              # 汇编器 as：把汇编翻译成二进制的 sum.o
$ gcc sum.c -o sum          # 链接器 ld：把 sum.o 和库拼成可执行文件

$ ./sum
15

$ file sum.o
sum.o: ELF 64-bit LSB relocatable, x86-64, version 1 (SYSV), not stripped
$ file sum
sum:   ELF 64-bit LSB pie executable, x86-64, dynamically linked, ...
\end{lstlisting}

sum.i 打开是展开后的 C 代码，通常有上千行；sum.s 就是汇编。而 sum.o 和 sum 用文本编辑器打开只会是乱码——file 命令一眼就能把它们认出来，这正是上图中蓝色和绿色的分界。

\subsubsection{反汇编：机器码与汇编的对应}

.o 和可执行文件里装的是二进制指令，要用 objdump 才能看。\texttt{objdump -d} 会把每一段机器码翻译回汇编，并把原始字节一并列出来：

\begin{lstlisting}[style=MyCppStyle, language=sh]
$ gcc -O1 -c sum.c
$ objdump -d sum.o
\end{lstlisting}

\begin{asmcode}
0000000000000000 <sum>:
   0:   f3 0f 1e fa           endbr64
   4:   48 85 f6              test   %rsi,%rsi
   7:   7e 27                 jle    30 <sum+0x30>
   9:   48 89 f8              mov    %rdi,%rax
   c:   48 8d 0c f7           lea    (%rdi,%rsi,8),%rcx
  10:   ba 00 00 00 00        mov    $0x0,%edx
  20:   48 03 10              add    (%rax),%rdx
  23:   48 83 c0 08           add    $0x8,%rax
  27:   48 39 c8              cmp    %rcx,%rax
  2a:   75 f4                 jne    20 <sum+0x20>
  2c:   48 89 d0              mov    %rdx,%rax
  2f:   c3                    ret
  30:   ba 00 00 00 00        mov    $0x0,%edx
  35:   eb f5                 jmp    2c <sum+0x2c>
\end{asmcode}

最左边是地址，中间是机器码字节，右边是 objdump 还原出的汇编：一条汇编恰好对应中间那串字节。带立即数的 mov 占 5 字节，不带立即数的 mov 和 add 占 3 字节，ret 只占 1 字节。这段反汇编与后文 \texttt{gcc -S} 生成的 sum.s 逐行一致，只是多了地址和机器码两列——所以读汇编和读机器码是同一件事，区别只在于后者要查指令编码表。

（上面省略了 gcc 为对齐循环入口插入的 8 字节 \texttt{nop} 填充，所以地址从 0x10 直接跳到 0x20。）

需要注意的是，objdump 默认也用 AT\&T 语法，但它不写 \texttt{movq}、\texttt{addq} 这类宽度后缀——宽度已经由 %rsi、%rdx 这些寄存器名确定了。

\subsubsection{汇编语言的特点}

先看一段 C 代码，以及 gcc -O1 为它生成的汇编（省略 gcc 自动加上、与语义无关的 \texttt{.cfi} 伪指令和 \texttt{endbr64}）：

\begin{lstlisting}[style=MyCppStyle]
long sum(long *a, long n) {
    long s = 0;
    for (long i = 0; i < n; i++)
        s += a[i];
    return s;
}
\end{lstlisting}

\begin{asmcode}
sum:
    testq   %rsi, %rsi          # n <= 0 ?
    jle     .L4
    movq    %rdi, %rax          # %rax = a
    leaq    (%rdi,%rsi,8), %rcx # %rcx = a + n*8，指向数组末尾
    movl    $0, %edx            # s = 0
.L3:
    addq    (%rax), %rdx        # s += *a
    addq    $8, %rax            # a++
    cmpq    %rcx, %rax          # 走到末尾了吗
    jne     .L3
.L1:
    movq    %rdx, %rax          # 返回值必须放进 %rax
    ret
.L4:
    movl    $0, %edx            # n <= 0 时 s = 0
    jmp     .L1
\end{asmcode}

\paragraph{读取}
汇编不提供结构化的语义，读它只能沿着 PC 的顺序走：先看助记符确定“做什么”，再看操作数确定“对谁做”，遇到 \texttt{jmp}、\texttt{jne} 这类跳转就把控制流接到目标标签上。所以读汇编的基本功不是逐行理解，而是跟着跳转指令画出控制流图——上面这段代码真正的循环体只有 \texttt{.L3} 到 \texttt{jne} 这四行，前面的 \texttt{testq}、\texttt{jle} 只是进入循环前的边界检查。

\paragraph{逻辑}
汇编里没有变量、类型和表达式，只有寄存器、内存和立即数三种东西。任何中间结果都必须显式占用一个具体位置：要么留在寄存器里，要么压到栈上（关掉优化后，gcc 会把每个 C 变量都摊在栈帧里）。运算只能在寄存器之间进行，内存中的数据要先 load 进寄存器、算完再 store 回去。控制流也不是 \texttt{if}、\texttt{for} 这样的结构，而是拆成“比较 + 跳转”两步：\texttt{testq}、\texttt{cmpq} 只负责设置条件码（Condition codes），\texttt{jle}、\texttt{jne} 再根据条件码决定跳不跳——这正是前面 CPU 模型里 Condition codes 存在的意义。

\paragraph{表示特点}
\begin{itemize}
    \item \textbf{一条汇编指令对应一条机器指令}：助记符就是操作码的文本写法，二者一一对应。因此汇编与 ISA 强绑定，同一段 C 在 x86-64 和 ARM 上生成的汇编完全不同，汇编层面不存在跨平台的抽象。
    \item \textbf{宽度写在名字里}：\texttt{movb}、\texttt{movw}、\texttt{movl}、\texttt{movq} 分别是 1、2、4、8 字节的操作；寄存器名同样自带宽度，\texttt{\%al} $\subset$ \texttt{\%ax} $\subset$ \texttt{\%eax} $\subset$ \texttt{\%rax}。
    \item \textbf{前缀区分操作数种类}：寄存器加 \texttt{\%}，立即数加 \texttt{\$}，操作数顺序是“源在前、目的在后”。这是 AT\&T 语法；Intel 语法没有前缀、顺序相反，两者只是同一台机器上的两种写法。
    \item \textbf{访存统一写成 \texttt{D(Rb,Ri,S)}}，即 地址 $=$ 偏移 $+$ 基址 $+$ 变址 $\times$ 比例。例子中的 \texttt{(\%rdi,\%rsi,8)} 就是数组下标寻址：基址是数组首地址，变址是下标，比例是 \texttt{sizeof(long)}。
    \item \textbf{参数和返回值的位置由调用约定固定}：前六个整型参数依次放在 \texttt{\%rdi}、\texttt{\%rsi}、\texttt{\%rdx}、\texttt{\%rcx}、\texttt{\%r8}、\texttt{\%r9}，返回值放在 \texttt{\%rax}。所以 \texttt{sum} 一进来，\texttt{a} 已经在 \texttt{\%rdi}、\texttt{n} 已经在 \texttt{\%rsi} 里了。
\end{itemize}

\subsubsection{汇编操作}

\paragraph{寄存器}

寄存器在本课程的系统中

\paragraph{Move指令}

Move指令的操作数有三种：寄存器，内存地址和立即数。除了不能用一条指令在两个内存地址之间传递信息，其余可以两两互相传递。想要实现前一句的效果，应当经过寄存器中转。前面的寄存器表示源，后面的寄存器表示目标。

\begin{asmcode}
    movq %rax, %rdx
\end{asmcode}

\paragraph{内存寻址}

内存中寻找地址有两种寻址方法：一是直接从寄存器中取出一个数据作为内存地址，从寄存器中寻址；第二种是通用的寻址形式。

$$D(Rb,Ri,S)$$

Rb叫作基址，Ri是变址，S为比例因子（只允许1,2,4,8），D为常数偏移量。

$$Mem[Reg[Rb]+S*Reg[Ri]+D]$$

上面是寻址公式，可以覆盖大量的不同场景的寻址方式和计算方法。根据以上公式，地址计算的汇编指令是：

\begin{asmcode}
    leaq Src Dst
\end{asmcode}

这种汇编指令主要用于以下两种场景：

\begin{itemize}
    \item 在不访问的情况下计算内存地址
    \item 进行形如$x + k * y$的算术运算
\end{itemize}

譬如以下计算x*12的代码,利用寻址算法减少了内存寻址和计算开销。
\begin{asmcode}
    leaq (%rdi,%rdi,2), %rax 
    salq $2, %rax 
\end{asmcode}

\paragraph{算术运算指令和逻辑运算指令}

首先是双操作数的运算指令，写法是助记符后面跟 \texttt{S, D}，效果是拿 D 和 S 运算后写回 D：

\begin{itemize}
    \item \texttt{add S, D}：$D \leftarrow D + S$，加法
    \item \texttt{sub S, D}：$D \leftarrow D - S$，减法
    \item \texttt{imul S, D}：$D \leftarrow D \times S$，有符号乘法
    \item \texttt{and S, D}：$D \leftarrow D \wedge S$，按位与
    \item \texttt{or S, D}：$D \leftarrow D \vee S$，按位或
    \item \texttt{xor S, D}：$D \leftarrow D \oplus S$，按位异或
    \item \texttt{sal k, D}、\texttt{shl k, D}：$D \leftarrow D \ll k$，左移，两者完全等价
    \item \texttt{sar k, D}：算术右移 $k$ 位，空位补符号位
    \item \texttt{shr k, D}：逻辑右移 $k$ 位，空位补 $0$
\end{itemize}

其中 S 可以是立即数、寄存器或内存，D 可以是寄存器或内存，但 S 和 D 不能同时是内存。除移位外，其余指令都会顺带设置条件码。

接下来是单操作数的运算指令，操作数既是源又是目的：

\begin{itemize}
    \item \texttt{inc D}：$D \leftarrow D + 1$，自增；注意它\textbf{不}影响进位标志 CF
    \item \texttt{dec D}：$D \leftarrow D - 1$，自减；同样不影响 CF
    \item \texttt{neg D}：$D \leftarrow -D$，取负
    \item \texttt{not D}：$D \leftarrow \neg D$，按位取反；它\textbf{不}设置任何条件码
\end{itemize}

\end{document}